\begin{derivation}[储库真空态中的电场关联]
把正频场写成

\begin{equation*}
\hat E_i^{(+)}(\mathbf r,\omega)
=\int\dd^3 s\,K_{ia}(\mathbf r,\mathbf s;\omega)
\hat f_a(\mathbf s,\omega),
\end{equation*}
其中 \(K\) 包含正文\eqnrefs{eq:macqed-Eplus-SI-r12} 中的全部数值核。真空满足
\begin{equation*}
\langle0|\hat f_a(\mathbf s,\omega)
\hat f_b^\dagger(\mathbf s',\omega')|0\rangle
=\delta_{ab}\delta^{(3)}(\mathbf s-\mathbf s')\delta(\omega-\omega').
\end{equation*}
代入两份场算符后，储库的空间 \(\delta\) 函数消去一个积分：
\begin{align*}
&\langle0|\hat{\mathbf E}^{(+)}(\mathbf r,\omega)
\hat{\mathbf E}^{(-)}(\mathbf r',\omega')|0\rangle\\
&=\mu_0^2\omega^4\frac{\hbar\varepsilon_0}{\pi}
\delta(\omega-\omega')
\int\dd^3 s\,
\mathbf G^R(\mathbf r,\mathbf s;\omega)
\cdot\operatorname{Im}\boldsymbol\varepsilon_r(\mathbf s,\omega)
\cdot(\mathbf G^R)^\dagger(\mathbf r',\mathbf s;\omega).
\end{align*}
利用 \(\mu_0\varepsilon_0=1/c^2\) 和正文\eqnrefs{eq:green-optical-identity-SI-r12}，积分等于 \((c^2/\omega^2)\operatorname{Im}\mathbf G^R\)。于是
\begin{equation*}
\langle0|\hat{\mathbf E}^{(+)}(\mathbf r,\omega)
\hat{\mathbf E}^{(-)}(\mathbf r',\omega')|0\rangle
=\frac{\hbar\mu_0}{\pi}\omega^2
\operatorname{Im}\mathbf G^R(\mathbf r,\mathbf r';\omega)
\delta(\omega-\omega'),
\end{equation*}
即正文\eqnrefs{eq:macqed-vacuum-field-correlation-dp68}。
\end{derivation}

\begin{derivation}[电磁顶角在模式基中的矩阵元]
正能电子场与其 Dirac 伴随在固定时刻分别写成
\[
\hat\psi_+(\mathbf r)
=\sum_s\int\frac{\dd^3p}{(2\pi\hbar)^{3/2}}
\sqrt{\frac{c}{2E_{\mathbf p}}}
\,u_s(p)\hat c_{\mathbf ps}
 \ee^{\ii\mathbf p\cdot\mathbf r/\hbar},
\]
\[
\hat{\bar\psi}_+(\mathbf r)
=\sum_{s'}\int\frac{\dd^3p'}{(2\pi\hbar)^{3/2}}
\sqrt{\frac{c}{2E_{\mathbf p'}}}
\,\hat c_{\mathbf p's'}^\dagger\bar u_{s'}(p')
 \ee^{-\ii\mathbf p'\cdot\mathbf r/\hbar}.
\]
取矢势正频部分
\[
\hat{\mathbf A}^{(+)}_T(\mathbf r)
=\sum_\lambda\mathbf A_\lambda(\mathbf r)\hat a_\lambda.
\]
代入
\[
\hat H_{jA}=-q_{\mathrm e} c\int\dd^3r:
\hat{\bar\psi}\,\boldsymbol\gamma\cdot\hat{\mathbf A}_T\hat\psi:
\]
并抽取算符单项
$\hat c_{\mathbf p's'}^\dagger\hat c_{\mathbf ps}\hat a_\lambda$，其系数为
\begin{align*}
K^{(\lambda)}_{\mathbf p's',\mathbf ps}
={}&-q_{\mathrm e} c
\sqrt{\frac{c}{2E_{\mathbf p'}}}
\sqrt{\frac{c}{2E_{\mathbf p}}}\\
&\times\int\frac{\dd^3r}{(2\pi\hbar)^3}
\bar u_{s'}(p')
\boldsymbol\gamma\cdot\mathbf A_\lambda(\mathbf r)
 u_s(p)
 \ee^{\ii(\mathbf p-\mathbf p')\cdot\mathbf r/\hbar}.
\end{align*}
为了证明这确实就是物理通道矩阵元，而不只是“比较系数”，用 Dirac-$\delta$ 归一化一电子态
\[
|\mathbf p,s;1_\lambda\rangle
\equiv
\hat c_{\mathbf ps}^\dagger\hat a_\lambda^\dagger|0\rangle
\]
计算
\begin{align*}
&\langle0|\hat a_{\lambda'}
\hat c_{\mathbf p_f s_f}
\hat H_{jA}^{(+)}
\hat c_{\mathbf p_i s_i}^\dagger
\hat a_\lambda^\dagger|0\rangle\\
&\qquad=
K^{(\lambda)}_{\mathbf p_f s_f,\mathbf p_i s_i}
\,\delta_{\lambda'\lambda}.
\end{align*}
这里电子部分使用
\[
\{\hat c_{\mathbf ps},\hat c_{\mathbf p's'}^\dagger\}
=\delta_{ss'}\delta^{(3)}(\mathbf p-\mathbf p')
\]
两次收缩，光子部分使用
$\langle0|\hat a_{\lambda'}\hat a_\lambda^\dagger|0\rangle=\delta_{\lambda'\lambda}$；所有没有与外态匹配的算符项都湮灭真空。因而有效哈密顿量 写成
\[
\hat H_{jA}^{(+)}
=\hbar\sum_\lambda\int\dd^3p\dd^3p'\,
\mathcal G^{(\lambda)}_{\mathbf p's',\mathbf ps}
\hat c_{\mathbf p's'}^\dagger\hat c_{\mathbf ps}\hat a_\lambda
\]
时，必须有
\[
\hbar\mathcal G^{(\lambda)}_{\mathbf p's',\mathbf ps}
=K^{(\lambda)}_{\mathbf p's',\mathbf ps},
\]
即正文\eqnrefs{eq:qed_to_mode_matrix_element_detailed}。这样 $(2\pi\hbar)^{-3}$、两个 $\sqrt{c/(2E)}$ 与一量子模式归一化都由场展开和外态收缩唯一固定，而不是后来人为补上的归一化因子。
\end{derivation}
